\section{Introduction}

The question is whether zero-shot time-series foundation models improve probabilistic forecasts of the daily
number of new stars a GitHub repository receives. These series are counts, most days are zero for most
repositories, activity arrives in bursts tied to releases and external attention, and the panel one can
observe is survivor-biased: a repository that was deleted or made private has no history to fetch. Public
forecasting benchmarks are mostly smooth, seasonal, positive series, so this kind of data is a hard case for
zero-shot forecasting. The practical question is whether such a model is worth running at all on a sparse
count series, or whether a trailing mean does as well.

The story of this report is in four parts. A first version of the study, built in June and July 2026, found
two foundation models ahead of classical baselines on mean CRPS in every cell of its design. An adversarial
audit on 2026-07-12 found no code defect that invalidated the aggregate comparison, but two design flaws that
could not be fixed by rewriting prose: the cohorts had been selected on outcomes not available at forecast
time, and the evaluation windows mostly tested post-peak decay. We registered a redesign on 2026-09-12, before
any target was fetched, and ran it on 2026-09-13 without change. The revised result is less flattering and is
reported with the same care: the models win on the mean against a 28-day local mean, lose on most individual
repositories, and win on both summaries against a weekly AutoETS whose intervals are far too wide. A post-hoc
accounting of the saved run, labeled as such, says where the mean gain comes from.

This is not the first observation that an aggregate gain can coexist with a minority win rate. A July 2026
quick-commerce study found TimesFM's pooled gains concentrated in high-turnover items while it won a minority
of item-rounds on lumpy demand \citep{kim2026}, and a marketplace-panel study showed the aggregate and
per-series rankings of five zero-shot foundation models reversing on M5 \citep{islam2026}. Frozen foundation
models have also been evaluated on sparse daily counts before: twelve of them, the three used here among
them, on five sparse panels as an event-ranking task without a proper score \citep{schoess2026}; and TimesFM 3
on simulated intermittent demand, where its nine-decile pinball loss was no better than the history's
empirical deciles \citep{ebrahim2026}, a finding consistent with the typical-repository result reported here.
Earlier evaluations on cloud workloads found simple baselines ahead \citep{toner2025}, and the Chronos and
TimesFM papers themselves report benchmark results that include count-like series without a zero-day
diagnosis \citep{ansari2024, das2023}. Benchmark practice has moved toward reporting win rates alongside
relative errors \citep{fevbench2026}, and the forecast-evaluation literature warns against ratio metrics that
silently drop series \citep{hewamalage2023}. Star-count prediction on GitHub has been studied as a regression
on popular repositories \citep{borges2016}, and a live benchmark scores daily new stargazers on the 400
most-starred repositories \citep{garza2026}; the present design samples from a population rather than from
the popular. What this report adds is a registered, population-sampled, repository-clustered probabilistic
evaluation whose every number traces to a saved file.
